\documentclass[10pt,letterpaper]{article}
\usepackage[margin=0.76in]{geometry}
\usepackage{amsmath,amssymb,lmodern,microtype,enumitem,hyperref}
\hypersetup{colorlinks=true,urlcolor=blue!55!black}
\usepackage{xcolor}
\pagestyle{empty}
\setlength{\parindent}{0pt}
\setlength{\parskip}{5pt}
\begin{document}
{\LARGE\bfseries What pair correlations forget about moving zeros}\par
{\small Discussion note for Terence Tao \hfill 5 November 2026}

In your 2019 ACUE construction, the gaps are multiples of half the mean spacing, yet a whole range of low moments agree exactly with Haar unitary eigenvalues. This gave a finite model for asking how much the known statistics of zeta zeros can tell us about their spacing. We asked a related inverse question. Suppose one knows \emph{every} pair distance in a configuration. Can one tell how long its zeros survive under heat flow before two meet?

\textbf{Four points suffice.} Number eight equally spaced sites around a circle. Occupy
\[
 S=\{0,1,3,4\}\qquad\text{or its complement}\qquad C=\{2,5,6,7\}.
\]
Both have the same six shorter-arc distances, $1,1,2,3,3,4$. In fact, at every displacement $r$, the number of occupied pairs equals the number of vacant pairs. Expanding $(1-1_S(x))(1-1_S(x+r))$ and summing proves this for any half-filled circle.

Rotate both configurations by $-\pi/2$. Their points are the zeros of
\[
 F(x)=\cos 2x-\sqrt2\cos x+1,\qquad G(x)=\sin 2x+\sqrt2\sin x .
\]
Run ordinary scalar heat for time $s$, so mode $k$ decays by $e^{-k^2s}$. For $G$, the two moving roots satisfy $\cos x=-e^{3s}/\sqrt2$ and meet the fixed root at $\pi$ when $s=(\log 2)/6\approx0.11552$. For $F$, the quadratic in $\cos x$ first acquires a double root when $y=e^{-2s}$ reaches the largest root of $4y^3-4y+1=0$, at $s\approx0.08863$. Pair distances lose track of which neighbors pull on the same zero. Replacing $x$ by $rx$ gives the same distinction for $4r$ points, with both times divided by $r^2$.

\textbf{What extra information matters?} On twelve sites take $A=\{0,1,2,4,7,8\}$ and its complement. Even their sorted consecutive gaps agree, $(1,1,1,2,3,4)$. Write $p_k=\sum_j e^{ik\theta_j}$. The cubic $\mathrm{Re}(p_1^2\overline{p_2})$ takes values $2$ and $-2$, giving different initial slopes of $|p_2|^2$ under heat. Moving a little ACUE probability from one symmetry orbit to the other therefore preserves the \emph{joint law} of the full pair data and gap histogram while changing the mean heat response.

\textbf{Where the distinction reaches.} The first collision for Haar points occurs on the $N^{-8/3}$ scale, while ACUE's lattice gives a lower bound $\pi^2/(8N^2)$. Local convergence alone cannot give the Haar result. Moving one uniformly chosen point next to another can force an arbitrarily small gap while changing the law in a fixed unfolded window of length $L$ by at most $2L/N$ in total variation. Under circular Dyson Brownian motion, multitime trace moments agree for ACUE and CUE up to total weight $N$ in each sign, and an explicit witness separates them at $N+1$. Our six-point perturbation already changes a cubic inside that protected range. For zeta, the open step is an arithmetic estimate beyond it.

\textbf{Five questions for discussion}
\begin{enumerate}[leftmargin=*,nosep]
\item Which three-zero statistic would you try first beyond the protected range, after subtracting all contributions already determined by pairs?
\item Can we perturb the ACUE weights while preserving the protected mixed moments as well as the pair data, and still change the first collision time?
\item What is the weakest estimate on the neighbors of one close pair that would let us predict its collision time for zeta?
\item Under RH, our frequency-two test reduces to a centered variance of short prime sums. Which signed term near $X=T^2$ offers the best chance of a strict saving?
\item Which family of function-field $L$-functions would let us compute a first unprotected statistic directly from primes and compare it with both ensembles?
\end{enumerate}
{\small \href{https://terrytao.wordpress.com/2019/05/08/the-alternative-hypothesis-for-unitary-matrices/}{Tao's ACUE note (2019)}\quad
\href{https://github.com/QingyunSun/Riemann-hypothesis-and-random-matrix/tree/4b96cd9/research/counterexamples}{Constructions, proofs and exact checks}}
\end{document}
